\documentclass[xcolor=dvipsnames, compress, 12pt, aspectratio=169, handout]{beamer}
\usepackage{../preamble_slides}

\title{H-1B Visa Project: Updates}
\author{Amy Kim}
\date{August 13, 2025}

\begin{document}
\begin{frame}[plain]
    \titlepage
    \addtocounter{framenumber}{-1}
\end{frame}

\begin{frame}
    \frametitle{Updates}
    \begin{enumerate}
        \item \textbf{Benchmarking Regression Results}
        \begin{itemize}
            \item Last time: regression results where outcome indicator for still being at company in $t+1, t+2,...$ 
            \item What would we actually expect the coefficient to be with no measurement error? How do my results compare?
        \end{itemize}
        \vspace{2mm}
        \item \textbf{New Results}
        \begin{itemize}
            \item What is actually happening to winners/losers? (Staying in U.S., return to school, promotions, etc.)
        \end{itemize}
        \vspace{2mm}
        \item \textbf{Miscellaneous}
        \begin{itemize}
            \item Updates on questions from last time and potential data issues
        \end{itemize}
    \end{enumerate}
\end{frame}

\section{Benchmarking Regression Results}
\begin{frame}
    \frametitle{Empirical Strategy: Basic Version}
For H-1B lottery application $j$:
\begin{equation*}
    stay^{1yr}_{i(j), c(j), t(j)} = \beta_0 + \beta_1 win_{j} + \varepsilon{j}
\end{equation*}

where $i(j)$, $c(j)$, and $t(j)$ denote the individual, company, and lottery year respectively of application $j$, $stay_{i,c,t}$ indicates whether individual $i$ who was observed working at company $c$ in year $t$ is still working at the firm one year later (in $t+1$). $win_{j}$ indicates whether application $j$ was selected in the H-1B lottery.

\vspace{2mm}
Our coefficient of interest is $\beta_1$, the effect of winning the lottery on staying at the firm (suppressing some notation for convenience):
\begin{equation*}
    \beta_1 = \mathbb{E}[stay_{ict}|win_j = 1] - \mathbb{E}[stay_{ict}|win_j = 0]
\end{equation*} 
\vspace{2mm}

In a world where all lottery losers leave immediately ($\mathbb{E}[stay_{ict}|win_j = 0] = 0$) and everyone else stays with probability $r$ ($\mathbb{E}[stay_{ict}|win_j = 1] = r$), we have $\beta_1 = r$. 

\end{frame}

\begin{frame}
\frametitle{Complication 1: Losers Don't Leave Immediately}
In my setting, the first reason why I might be estimating a smaller $\beta_1$ than expected is because of the repeated draw structure of the lottery.
\vspace{2mm}

Let's assume that all applicants $i(j)$ are on 3-year OPT visas, so they apply for the H-1B visa every year until they win, with a max of three attempts. Applicants who are on their third try and lose leave the company immediately, otherwise they stay with probability $r$. Then:
\begin{align*}
    \mathbb{E}[stay_{ict}|win_j = 0] &= \mathbb{E}[\mathbb{E}[stay_{ict}|win_j=0, attempt(j) = a] | win_j = 0] \\
    &= \mathbb{P}[attempt(j) < 3] \cdot r + \mathbb{P}[attempt(j) = 3] \cdot 0 \\
    &= (1-s_3) \cdot r
\end{align*} 
where $attempt(j) \in \{1,2,3\}$ denotes the `attempt number' of application $j$ and $s_3$ represents the share of applicants that are on their third attempt.
\end{frame}

\begin{frame}
    \frametitle{Complication 2: Multiplicity/Mismeasurement}
    The second complication is when I measure $stay_{ict}$, I'm really measuring a weighted average of the true $stay_{ict}$ and the $stay_{i'ct}$ for some falsely matched individuals $i'$. If the true individual has weight $w$ on average and the falsely matched individuals have average stay rate $q$:

    \begin{align*}
        \mathbb{E}[\hat{\beta}_1] &= w\cdot r + (1-w) \cdot q - (w \cdot (1-s_3) \cdot r + (1-w) \cdot q) \\ 
        &= w s_3 r
    \end{align*}

    Note that all estimated coefficients (not just for outcome stay rate) will be downweighted by $w$ under these assumptions.
\end{frame}

\begin{frame}
    \frametitle{Putting it All Together}
    \begin{equation*}
        \mathbb{E}[\hat{\beta}_1] = w s_3 r
    \end{equation*}
    \begin{itemize}
        \item Estimating $r$ (average stay rate of winners): various methods yield similar estimates of 0.5
        \item Estimating $s_3$ (share of applicants in sample that are on their third try): assuming constant win probability of $0.4$ in sample yields estimate of 0.18
        \item Estimating $w$ (average weight on `true match'): depending on merge specification, estimates range from 0.5 to 0.7
    \end{itemize}
    \begin{align*}
        \hat{w} \cdot \hat{s}_3 \cdot \hat{r} \in [0.5 \cdot 0.18 \cdot 0.5, 0.7 \cdot 0.18 \cdot 0.5] = [0.045, 0.063]
    \end{align*}
    So even if I'm estimating and measuring everything correctly, in the best matched sample I would only expect to see a coefficient of $0.06$.
\end{frame}

\begin{frame}
    \frametitle{Some Implications}
    \begin{equation}
        \mathbb{E}[\hat{\beta}_1] = w s_3 r
    \end{equation}
    \begin{enumerate}
        \item The estimated coefficient on $stay_{ict}$ should increase as $s_3$ (the share of the sample that is on their third try) increases.
        \begin{itemize}
            \item I can vary this either by restricting the lottery year ($s_3$ should increase in later years as the lottery gets harder), or by measuring outcome $stay^{2yr}$ (the resulting coefficient would be downweighted instead by $s_2 > s_3$)
            \item If the estimated coefficient does not increase with $s_3$, it suggests that my ex-ante three year OPT model of behavior is incorrect (or that there is mismeasurement)
        \end{itemize}
        \item The estimated coefficient on $stay_{ict}$ should increase as $w$ (the average weight of the `true' individual) increases.
        \begin{itemize}
            \item I can vary this in the `postfiltered' merge sample, which is explicitly filtered on multiplicity
            \item If the estimated coefficient does not increase with $w$, it suggests mismeasurement
        \end{itemize}
    \end{enumerate}
\end{frame}


\begin{frame}
    \frametitle{Comparing Benchmarks to Real Results}
    Baseline regression + control for ADE and firm x year fixed effects:
    \small
    \centering
\resizebox{0.85\textwidth}{!}{
\begin{tabular}{lcccccc}
\toprule
& \multicolumn{2}{c}{\textbf{1. Low Multiplicity}} 
& \multicolumn{2}{c}{\textbf{2. Medium Multiplicity}} 
& \multicolumn{2}{c}{\textbf{3. High Multiplicity}} \\
\cmidrule(lr){2-3} \cmidrule(lr){4-5} \cmidrule(lr){6-7}
& stay1yr & stay2yr &  stay1yr & stay2yr & stay1yr & stay2yr \\
& (1) & (2) & (3) & (4) & (5) & (6) \\
\midrule
win & 0.0149   &   0.0318**  &    0.0155**   &   0.0171**    & 0.0154***  &    0.0126** \\
& (0.0124)   &   (0.0135)     & (0.0072)   &   (0.0078)    &  (0.0056)   &   (0.0061) \\
\midrule
\textbf{Obs.}      & 2920     &    2920     &     7007     &      7007      &      10363        &     10363 \\
\midrule
\textbf{DV mean for winners} ($\hat{r}$) & 0.507 & 0.420 & 0.506 & 0.417 & 0.507 & 0.416 \\
$\hat{s}$ & 0.18 & 0.3 & 0.18 & 0.3 & 0.18 & 0.3 \\
$\hat{w}$ & 0.790 & 0.790 & 0.620 & 0.620 & 0.536 & 0.536 \\
\textbf{Predicted coefficient} ($wsr$) & 0.072 & 0.100 & 0.057 & 0.078 & 0.049 & 0.067 \\
% Predicted relative coef size & 1 & 1.39 & 0.79 & 1.08 & 0.68 & 0.93 \\
% Actual relative coef size & 1 & 2.13 & 1.04 & 1.15 & 1.03 & 0.85 \\
\bottomrule
\end{tabular}
}

\end{frame}
\section{New Results}
\begin{frame}
    \frametitle{New Outcome Variables}
    \begin{itemize}
        \item us1yr and us2yr: indicator for working in US 1 and 2 yrs after lottery 
        \item promote1yr and promote2yr: indicator for working at same firm but switching jobs 1 and 2 yrs after lottery 
        \item school1yr and school2yr: indicator for attending school 1 and 2 yrs after lottery 
        \item Also have Revelio imputed salary, but it seems really bad so looking for a way to improve on that (maybe looking at indicators for different jobs)
    \end{itemize}
\end{frame}

\begin{frame}
    \frametitle{New Results}
    Baseline regression + control for ADE and firm x year fixed effects (medium multiplicity sample):
    \vspace{2mm}
    \small


    \centering
\resizebox{0.85\textwidth}{!}{
\begin{tabular}{lcccccc}
\toprule
& us1yr & us2yr &  promote1yr & promote2yr & school1yr & school2yr \\
& (1) & (2) & (3) & (4) & (5) & (6) \\
\midrule
win & 0.0079**  &      0.0077     &  -0.0015    &    0.0015   &     0.0030   &      0.0090 \\
& (0.0034)  &    (0.0047)  &    (0.0064)   &   (0.0072)   &   (0.0075)    &   (0.0076) \\
\midrule
\textbf{Obs.}      & 7007     &      7007     &     7007     &      7007      &     7007     &      7007  \\
\midrule
\textbf{DV mean} & 0.953 & 0.905 & 0.192 & 0.271 & 0.341 & 0.372 \\
\bottomrule
\end{tabular}
}

\end{frame}


\begin{frame}
    \frametitle{Miscellaneous Updates}
\begin{itemize}
            \item Evaluating merge quality for winners (more info on winner applications): seems to be pretty good, more than half of applications are matched to at least one indiv with exact match on field of study or job title (pretty restrictive) 
            \item From last time: winners and losers balanced on multiplicity, so any match issues shouldn't create bias
            \item Data quality issues: trying to contact Revelio rep through Charissa
            \item Previous result on updated date incorrect: I think updated date actually refers to date Revelio data was refreshed from LinkedIn, not date profile was updated
\end{itemize}
\end{frame}

\begin{frame}
    \frametitle{Next Steps}
    Potential paths forward 
    \begin{itemize}
        \item Go back and work on improving matching (putting time limit of a month or two)
        \item Try to start putting together working draft
        \item Merge in external data sources 
    \end{itemize}
\end{frame}

\appendix


\begin{frame}
    \frametitle{Estimating $w$ in Data}

    \begin{itemize}
        \item \textbf{Method 1:} If I assume the match with the highest weight is the `true match' then I can estimate $\hat{w} = \frac{1}{n} \sum_{j} \max_{i' \in \mathcal{I}(j)} w_{i'j}$ where $\mathcal{I}(j)$ is the set of individuals to which application $j$ is matched (average of max weight)
        \begin{itemize}
            \item Under this method I estimate $\hat{w} \in [0.55,0.7]$ for various merge specifications
        \end{itemize}
        \vspace{2mm}
        \item \textbf{Method 2:} If I instead assume the `true match' has an equal probability of being any of the potential matches regardless of weight, I can estimate $\hat{w} = \frac{1}{n} \sum_{j} \frac{1}{|\mathcal{I}(j)|}\sum_{i' \in \mathcal{I}(j)} w_{i'j}$ (average of average weight)\footnote{Last time I talked about multiplicity, the avg number of LinkedIn profiles to which a given application is matched, which ranges from 2 to 8 in these samples. Last time we talked about inverting this number to get w, which would imply 0.1 to 0.5. But I think averaging after inverting makes more sense when thinking about average weight of the true individual.}
        \begin{itemize}
            \item Under this method I estimate $\hat{w} \in [0.5,0.65]$ for various merge specifications
        \end{itemize}
    \end{itemize}
\end{frame}


\begin{frame}
    \frametitle{Estimating $s_3$ in Data}
    In reality $s_3$ is lottery year dependent, but for now assume win probability $p$ is constant across years and the number of applicants is also constant across years (not far from true in our sample). Then:
    \begin{equation*}
        s_3 = \frac{(1-p)^2}{1 + (1-p) + (1-p)^2}
    \end{equation*}
    In our merged sample I estimate $\hat{p} \approx 0.4$, so $\hat{s}_3 \approx 0.18$. In reality since $p$ is actually decreasing over time, $s_3$ is larger in later lottery years.
\end{frame}

\end{document}
